% ============================================================================
\section{Análisis competitivo}

El análisis comparó Bloomberg, Power BI, Pyramid Analytics, ThoughtSpot, Collibra y la práctica
actual de consultar fuentes separadas. Las alternativas financieras ofrecen profundidad de
mercado; las herramientas de inteligencia de negocio facilitan visualización y autoservicio; y
las plataformas de gobierno organizan catálogo y linaje. Ninguna reunió en una misma propuesta
catálogo gobernado, trazabilidad visible, consulta en lenguaje natural y acceso institucional sin
un costo asociado al número de usuarios.

\figuraux{figuras/a3_matriz_competitiva.png}
{Matriz competitiva utilizada para situar la propuesta de Karisma Data.}{fig:a5-competencia}

El posicionamiento elegido no intenta reemplazar todas las capacidades especializadas. Karisma
Data funciona como puerta institucional: ayuda a descubrir, comprender y llevar un dato a la
herramienta adecuada, conservando permisos y procedencia.

% ============================================================================
\section{Card sorting}

La Actividad 3 utilizó 35 tarjetas para contrastar el vocabulario y las agrupaciones propuestas.
El ejercicio abierto produjo nueve conglomerados, mientras que los dominios de negocio mostraron
el mayor acuerdo. La principal presión apareció en elementos transversales, especialmente consumo
por API y calidad de datos, que no debían forzarse a una sola categoría.

\begin{uxtabla}{L{4.0cm}Y}{Resultados del card sorting y decisión aplicada}
  \uxheadrow \thd{Resultado} & \thd{Decisión de arquitectura} \\
  Acuerdo alto en crédito, liquidez, derivados y otros dominios & Mantener una entrada temática reconocible dentro de Exploración y extracción \\
  Nueve conglomerados frente a cuatro categorías de primer nivel & Resolver parte de la diversidad en segundo nivel y no multiplicar la navegación principal \\
  Consumo por API y calidad solicitados en más de un grupo & Usar facetas y accesos cruzados, sin duplicar el contenido \\
  Indicadores de riesgo buscados en exploración & Mover tableros e indicadores desde Gobierno del dato a Exploración y extracción \\
\end{uxtabla}

\figuraux{figuras/a3_dendrograma_a5.png}
{Agrupamiento resultante del ejercicio de card sorting.}{fig:a5-card-sorting}

El ejercicio se realizó con ocho participantes y permitió comparar hipótesis de estructura antes
de construir las pantallas. Sus resultados orientaron la arquitectura que posteriormente se
evaluó mediante tareas observables en el demo.

% ============================================================================
\section{Arquitectura de información}

La arquitectura final organiza el portal en cuatro entradas: Inicio, Exploración y extracción,
Gobierno del dato y Administración. El asistente es transversal porque puede acompañar tareas de
varias categorías. El contenido se adapta por permisos, pero conserva nombres y rutas estables.

\figuraux{figuras/a3_arquitectura_informacion_a5.png}
{Arquitectura de información derivada del análisis y el card sorting.}{fig:a5-arquitectura}

\section{Mapa de navegación}

El mapa traduce la arquitectura a recorridos. Inicio reúne buscador, recientes, favoritos y
alertas; Exploración contiene catálogo, filtros, exportaciones y tableros; Gobierno concentra
definiciones y versiones; Administración cubre cuentas, aprobaciones, bitácora e integraciones.

\figuraux{figuras/a3_mapa_sitio.png}
{Mapa de navegación de Karisma Data. Los accesos transversales conectan funciones que sirven a más de una categoría.}{fig:a5-mapa}

Esta estructura mantiene revelación progresiva: la navegación principal ofrece pocas decisiones y
cada módulo abre las herramientas avanzadas solo cuando la tarea las requiere.
